% === Chapter 1: Introduction / How to cite & reference ======================
% このファイルは main.tex から \subfile{chapters/chap1} として読み込まれます。
% 単体ビルドも可：\documentclass[../main.tex]{subfiles}
\documentclass[../main.tex]{subfiles} % （subfiles パッケージを使っているので、親ファイル ../main.tex を指定します）
\begin{document}

\section{文献の引用}
% =====================================================================
% 1. 文献の引用（\cite{...}）
%  - 本テンプレでは citestyle=styles/thesis-cite により、\cite は上付き数字。
%  - 例："... と報告されている\cite{cite:Lundstrom2018}．"
%  - 和文では「句読点（．，）」の直前に置くと見栄えが良い。
% =====================================================================
フランシスの論文を引用する\cite{cite:Francis2023}．
奥山先生の論文を引用する\cite{cite:Okuyama2022}．
複数同時に引用する場合は自動で圧縮される\cite{cite:Francis2023, cite:Takanari2024,cite:Koo2016, cite:Bakdash2017}．

1行空けると改行する。文量稼ぎのために、研究室HPの研究紹介文を差し込みます。 近年、メカトロニクス技術の発展には目覚しいものがあるが、機械の自動化あるいは人間と機械との協調動作をさらに促進するためには、新しいセンサやアクチュエータの創製が必要である。このためには、例えば従来の材料には見られない機能性材料等の特性を効率良く引き出し、利用することを考えなければならない。機能性材料は今日まで種々提案されている一方で、その応用技術は今だ十分に確立していない。機能性材料のセンサ・アクチュエータとしての応用はこれからの問題であり、その研究は十分な発展性を持つ。

新しいセンサ・アクチュエータの開発は工学の重要な課題であるが、機能性材料のセンサ・アクチュエータへの応用、特に医学への応用問題は人口の高齢化を迎える現代社会において重要な課題である。機能性材料の設計と創製は材料学者が担う一方で、その応用技術の開発はおもに機械工学者が担当する分野である。材料の利用にあたっては両者の共同研究が必要であり、さらに医学への応用を考えるならば医学者の参画もなければならない。機能性材料の応用は学際問題であり、工学者と医学者の理解・共同研究に貢献するものである。

本研究室は計測制御に関する教育と研究を行っており、現在インテリジェントロボットハンドの開発、センサやアクチュエータの医学分野への適用、それらに計測・自動制御の組み入れたマイクロ･マクロメカトロニクスの研究を行っている。

具体的には、小型ロボット指の機構･駆動制御法の開発、人間の指感覚を備えたマルチスマート柔軟触感センサの開発、前立腺癌・乳がんの触診用プローブの開発、管腔臓器の人工開閉弁の開発に関する研究等を行っている。


\section{図の参照}
% =====================================================================
% 2. 図の参照（\figref{...}）
%  - マクロ \figref により "Fig. 1" の体裁で出力される。
%  - 図の画像は figures/chap1/dummy.pdf を想定。
%  - \graphicspath{\subfix{figures/}} 設定で親/子どちらのビルドでも解決。
% =====================================================================
\begin{figure}[bt]
  \centering
  % 図ファイルは figures/chap1/dummy.pdf を置いておくこと
  \includegraphics[width=0.6\linewidth]{chap1/dummy.pdf}
  \caption{Dummy illustration for citation examples.}
  \label{fig:dummy}
\end{figure}

本文では \figref{fig:dummy} のように記述する．
図番号の前後はマクロが自動で整えるため，単にラベルを渡せばよい．
\LaTeX 上の記述では，\verb|\begin{figure}[bt]| のように [bt] という配置指定を用いている．
ここで[b] は下端（bottom）， [t] はページ上端（top）への配置を許可する指定であり，
LaTeX はページ全体のバランスを考慮して最適な位置に自動で図を配置する．
必要に応じて [p]（図表のみのページ）や [h]（その位置への配置を希望），[!tbp]（配置制約を緩める）なども利用できる．
一般的には [tbp] あるいは [!tbp] を用いると，図が極端に離れたり詰まりすぎたりするのを防げる．

文量稼ぎのために、研究室HPの研究紹介文を差し込みます。 近年、メカトロニクス技術の発展には目覚しいものがあるが、機械の自動化あるいは人間と機械との協調動作をさらに促進するためには、新しいセンサやアクチュエータの創製が必要である。このためには、例えば従来の材料には見られない機能性材料等の特性を効率良く引き出し、利用することを考えなければならない。機能性材料は今日まで種々提案されている一方で、その応用技術は今だ十分に確立していない。機能性材料のセンサ・アクチュエータとしての応用はこれからの問題であり、その研究は十分な発展性を持つ。

新しいセンサ・アクチュエータの開発は工学の重要な課題であるが、機能性材料のセンサ・アクチュエータへの応用、特に医学への応用問題は人口の高齢化を迎える現代社会において重要な課題である。機能性材料の設計と創製は材料学者が担う一方で、その応用技術の開発はおもに機械工学者が担当する分野である。材料の利用にあたっては両者の共同研究が必要であり、さらに医学への応用を考えるならば医学者の参画もなければならない。機能性材料の応用は学際問題であり、工学者と医学者の理解・共同研究に貢献するものである。

本研究室は計測制御に関する教育と研究を行っており、現在インテリジェントロボットハンドの開発、センサやアクチュエータの医学分野への適用、それらに計測・自動制御の組み入れたマイクロ･マクロメカトロニクスの研究を行っている。

具体的には、小型ロボット指の機構･駆動制御法の開発、人間の指感覚を備えたマルチスマート柔軟触感センサの開発、前立腺癌・乳がんの触診用プローブの開発、管腔臓器の人工開閉弁の開発に関する研究等を行っている。

\section{表の参照}
% =====================================================================
% 3. 表の参照（\tabref{...}）
%  - マクロ \tabref により "Table 1" の体裁で出力される。
%  - 表番号は \label を付与したキャプションの表に対して自動採番。
% =====================================================================
\begin{table}[bt]
  \centering
  \caption{Example parameters.}
  \label{tab:example}
  \begin{tabular}{ll}
    \hline
    Name & Value\\
    \hline
    Sampling rate & 100\,kHz\\
    Window (RMS)  & 0.375\,s\\
    Frequency     & 250\,Hz\\
    \hline
  \end{tabular}
\end{table}

本文では \tabref{tab:example} のように参照する．
表番号の採番は自動で行われる．

文量稼ぎのために、研究室HPの研究紹介文を差し込みます。 近年、メカトロニクス技術の発展には目覚しいものがあるが、機械の自動化あるいは人間と機械との協調動作をさらに促進するためには、新しいセンサやアクチュエータの創製が必要である。このためには、例えば従来の材料には見られない機能性材料等の特性を効率良く引き出し、利用することを考えなければならない。機能性材料は今日まで種々提案されている一方で、その応用技術は今だ十分に確立していない。機能性材料のセンサ・アクチュエータとしての応用はこれからの問題であり、その研究は十分な発展性を持つ。

新しいセンサ・アクチュエータの開発は工学の重要な課題であるが、機能性材料のセンサ・アクチュエータへの応用、特に医学への応用問題は人口の高齢化を迎える現代社会において重要な課題である。機能性材料の設計と創製は材料学者が担う一方で、その応用技術の開発はおもに機械工学者が担当する分野である。材料の利用にあたっては両者の共同研究が必要であり、さらに医学への応用を考えるならば医学者の参画もなければならない。機能性材料の応用は学際問題であり、工学者と医学者の理解・共同研究に貢献するものである。

本研究室は計測制御に関する教育と研究を行っており、現在インテリジェントロボットハンドの開発、センサやアクチュエータの医学分野への適用、それらに計測・自動制御の組み入れたマイクロ･マクロメカトロニクスの研究を行っている。

具体的には、小型ロボット指の機構･駆動制御法の開発、人間の指感覚を備えたマルチスマート柔軟触感センサの開発、前立腺癌・乳がんの触診用プローブの開発、管腔臓器の人工開閉弁の開発に関する研究等を行っている。



\section{式の参照}
% =====================================================================
% 4. 式の参照（\eqnref{...}）
%  - マクロ \eqnref により "式(1)" の体裁で出力される。
%  - \label は equation 環境の中に置く。
% =====================================================================
運動方程式は \eqnref{eq:newton} のように参照する．
\begin{equation}
  \label{eq:newton}
  F = m a
\end{equation}
式番号は本文全体で連番となる．

文量稼ぎのために、研究室HPの研究紹介文を差し込みます。 近年、メカトロニクス技術の発展には目覚しいものがあるが、機械の自動化あるいは人間と機械との協調動作をさらに促進するためには、新しいセンサやアクチュエータの創製が必要である。このためには、例えば従来の材料には見られない機能性材料等の特性を効率良く引き出し、利用することを考えなければならない。機能性材料は今日まで種々提案されている一方で、その応用技術は今だ十分に確立していない。機能性材料のセンサ・アクチュエータとしての応用はこれからの問題であり、その研究は十分な発展性を持つ。

新しいセンサ・アクチュエータの開発は工学の重要な課題であるが、機能性材料のセンサ・アクチュエータへの応用、特に医学への応用問題は人口の高齢化を迎える現代社会において重要な課題である。機能性材料の設計と創製は材料学者が担う一方で、その応用技術の開発はおもに機械工学者が担当する分野である。材料の利用にあたっては両者の共同研究が必要であり、さらに医学への応用を考えるならば医学者の参画もなければならない。機能性材料の応用は学際問題であり、工学者と医学者の理解・共同研究に貢献するものである。

本研究室は計測制御に関する教育と研究を行っており、現在インテリジェントロボットハンドの開発、センサやアクチュエータの医学分野への適用、それらに計測・自動制御の組み入れたマイクロ･マクロメカトロニクスの研究を行っている。

具体的には、小型ロボット指の機構･駆動制御法の開発、人間の指感覚を備えたマルチスマート柔軟触感センサの開発、前立腺癌・乳がんの触診用プローブの開発、管腔臓器の人工開閉弁の開発に関する研究等を行っている。



\section{章の参照}
% =====================================================================
% 5. 章の参照（\chapref{...}）
%  - 本テンプレでは和文スタイルのマクロを用意（\chapref → 「3章」）。
%  - 参照先の章に \label を付けておく必要がある。
%    例： main.tex 側の「Methods」直後に \label{chap:methods} を置く。
% =====================================================================
本論文の \chapref{chap:methods} に biblatex（BibTeX）を用いた参考文献出力方法を示す．
章への参照は本文の流れを案内するのに有効である．

文量稼ぎのために、研究室HPの研究紹介文を差し込みます。 近年、メカトロニクス技術の発展には目覚しいものがあるが、機械の自動化あるいは人間と機械との協調動作をさらに促進するためには、新しいセンサやアクチュエータの創製が必要である。このためには、例えば従来の材料には見られない機能性材料等の特性を効率良く引き出し、利用することを考えなければならない。機能性材料は今日まで種々提案されている一方で、その応用技術は今だ十分に確立していない。機能性材料のセンサ・アクチュエータとしての応用はこれからの問題であり、その研究は十分な発展性を持つ。

新しいセンサ・アクチュエータの開発は工学の重要な課題であるが、機能性材料のセンサ・アクチュエータへの応用、特に医学への応用問題は人口の高齢化を迎える現代社会において重要な課題である。機能性材料の設計と創製は材料学者が担う一方で、その応用技術の開発はおもに機械工学者が担当する分野である。材料の利用にあたっては両者の共同研究が必要であり、さらに医学への応用を考えるならば医学者の参画もなければならない。機能性材料の応用は学際問題であり、工学者と医学者の理解・共同研究に貢献するものである。

本研究室は計測制御に関する教育と研究を行っており、現在インテリジェントロボットハンドの開発、センサやアクチュエータの医学分野への適用、それらに計測・自動制御の組み入れたマイクロ･マクロメカトロニクスの研究を行っている。

具体的には、小型ロボット指の機構･駆動制御法の開発、人間の指感覚を備えたマルチスマート柔軟触感センサの開発、前立腺癌・乳がんの触診用プローブの開発、管腔臓器の人工開閉弁の開発に関する研究等を行っている。

\subsection{偶数ページの空白について}
本テンプレートではクラスオプション openright を有効にしており，章は必ず右（奇数）ページから始まる．この体裁を守るため，\LaTeX はページの偶奇を調整する目的で\verb|\cleardoublepage|を自動挿入し，必要に応じて左（偶数）ページが空白のまま挿入される．PDF閲覧では空白が気になる場合もあるが，これは体裁上の仕様でありエラーではない．添削時などで空白ページを入れないことが求められる場合は，openright を外す（または openany を用いる／oneside にする）ことで回避できる．

% --- 補足：章・節・付録の参照マクロ（和文スタイル） -----------------
% \chapref{chap:methods} → 「3章」
% \secref{sec:methods}   → 「2.3節」
% \appref{app:spec}      → 「付録A」
% ※ それぞれの参照先で \label{chap:...}, \label{sec:...}, \label{app:...} を付けておくこと。

% 参考：もし節を参照するなら次のように書く：
% 例）\secref{sec:methods}で示したように, ...
% 付録も同様に \appref{app:hardware} と書けば「付録A」のように出力される。

\end{document}